\documentclass[12pt,a4paper]{article}

% ------------------------------------------------
% PACOTES
% ------------------------------------------------
\usepackage[utf8]{inputenc}
\usepackage[T1]{fontenc}
\usepackage[brazil]{babel}
\usepackage{amsmath,amssymb}
\usepackage{geometry}
\usepackage{graphicx}
\usepackage{booktabs}
\usepackage{float}
\usepackage{array}
\usepackage{longtable}
\usepackage{listings}
\usepackage{xcolor}
\usepackage{hyperref}

\geometry{a4paper, margin=2.5cm}
\setlength{\parindent}{1.25cm}
\setlength{\parskip}{0.2cm}

\lstset{
    language=Python,
    basicstyle=\ttfamily\footnotesize,
    breaklines=true,
    frame=single,
    numbers=left,
    numberstyle=\tiny,
    showstringspaces=false
}

% ------------------------------------------------
% INÍCIO DO DOCUMENTO
% ------------------------------------------------
\begin{document}

% ------------------------------------------------
% CAPA
% ------------------------------------------------
\begin{titlepage}
    \begin{center}
        {\large FATEC-RP}\\[0.5cm]
        {\large Análise e Desenvolvimento de Sistemas}\\[0.5cm]
        {\large Programação Linear}\\[2.5cm]
        {\Large \textbf{TRABALHO FINAL}}\\[0.5cm]
        {\large \textbf{Questão 2: Localização de Centros
        de Distribuição}}\\[3cm]
        \textbf{Integrantes:}\\[0.3cm]
        Adrian Souza\\
        Heitor Lopes\\
        Victor Breno\\
        Yuri Ramos\\
        Isaac Leonardo\\[2cm]
        Professor: Junior Bonafim\\[3cm]
        Ribeirão-Preto\\
        2026
    \end{center}
\end{titlepage}

% ------------------------------------------------
% SUMÁRIO
% ------------------------------------------------
\tableofcontents
\newpage

% ------------------------------------------------
% 1. INTRODUÇÃO
% ------------------------------------------------
\section{Introdução}

O presente trabalho aborda um problema de localização de
centros de distribuição, cujo objetivo é determinar quais
localidades devem receber instalações de centros para atender
um conjunto de localidades consumidoras.

O problema considera um conjunto de localidades, cada uma
com coordenadas cartesianas e uma quantidade estimada de
viagens semanais. Nem todas as localidades estão habilitadas
a receber um centro de distribuição; esse subconjunto é
denotado por $C \subseteq I$.

A questão apresenta dois objetivos distintos:

\begin{itemize}
    \item \textbf{Item (a):} minimizar a maior distância entre
    uma localidade e o centro ao qual ela está associada.
    \item \textbf{Item (b):} maximizar a quantidade total de
    viagens semanais das localidades que estejam dentro de um
    raio de cobertura $R$ de pelo menos um centro.
\end{itemize}

Para resolver os dois itens, são construídos modelos de
programação linear inteira mista (PLIM), utilizando variáveis
binárias para representar as decisões de instalação,
atribuição e cobertura. Os modelos são resolvidos com o
solver HiGHS, disponível na biblioteca \texttt{SciPy}.

% ------------------------------------------------
% 2. DESCRIÇÃO DO PROBLEMA
% ------------------------------------------------
\section{Descrição do Problema}

O problema consiste em escolher exatamente $m$ localidades
para receber centros de distribuição, considerando um
conjunto de $n$ localidades existentes.

Cada localidade possui quatro informações:

\begin{itemize}
    \item coordenada horizontal $x_i$;
    \item coordenada vertical $y_i$;
    \item quantidade de viagens semanais $v_i$;
    \item indicador de elegibilidade $e_i \in \{0,1\}$.
\end{itemize}

O indicador $e_i$ assume valor 1 quando a localidade pode
receber um centro de distribuição e 0 caso contrário.

A distância entre duas localidades é calculada pela distância
euclidiana:
\begin{equation}
d_{ij} = \sqrt{(x_i-x_j)^2+(y_i-y_j)^2}.
\end{equation}

\subsection{Item (a): minimização da maior distância}

Deve-se escolher exatamente $m$ localidades e atribuir cada
localidade existente a um dos centros instalados, de modo a
minimizar a maior distância entre uma localidade e o centro
responsável pelo seu atendimento.

\subsection{Item (b): maximização das viagens cobertas}

Deve-se escolher exatamente $m$ centros, considerando um
raio de cobertura $R$. Uma localidade é considerada coberta
quando existe pelo menos um centro instalado a uma distância
menor ou igual a $R$ dela. O objetivo é maximizar a soma das
viagens semanais das localidades cobertas.

% ------------------------------------------------
% 3. FORMULAÇÃO MATEMÁTICA
% ------------------------------------------------
\section{Formulação Matemática}

\subsection{Conjuntos e parâmetros}

\begin{itemize}
    \item $I = \{1,2,\ldots,n\}$: conjunto de localidades;
    \item $C \subseteq I$: localidades elegíveis para receber
    centros (aquelas com $e_i = 1$);
    \item $n$: número total de localidades;
    \item $m$: número exato de centros a instalar;
    \item $d_{ij}$: distância euclidiana entre $i$ e $j$;
    \item $v_i$: viagens semanais da localidade $i$;
    \item $R$: raio de cobertura (item (b)).
\end{itemize}

% ------------------------------------------------
% 3.1 ITEM A
% ------------------------------------------------
\subsection{Modelo do item (a)}

\subsubsection*{Variáveis de decisão}

\begin{equation}
z_j = \begin{cases}
1, & \text{se um centro é instalado em } j \in C,\\
0, & \text{caso contrário;}
\end{cases}
\end{equation}

\begin{equation}
a_{ij} = \begin{cases}
1, & \text{se a localidade } i \text{ é atendida pelo centro } j,\\
0, & \text{caso contrário;}
\end{cases}
\end{equation}

\begin{equation}
D \ge 0 \text{ é a maior distância de atendimento.}
\end{equation}

\subsubsection*{Função objetivo}

\begin{equation}
\boxed{\min \; D}
\end{equation}

\subsubsection*{Restrições e justificativas}

\textbf{(R1) Número exato de centros.}
\begin{equation}
\sum_{j\in C} z_j = m.
\end{equation}
Justificativa: o enunciado exige exatamente $m$ centros.

\textbf{(R2) Cada cliente é atendido por exatamente um centro.}
\begin{equation}
\sum_{j\in C} a_{ij} = 1, \qquad \forall i\in I.
\end{equation}
Justificativa: atribuição única, exigida pelo item (a).

\textbf{(R3) Só é possível atribuir a um centro instalado.}
\begin{equation}
a_{ij} \le z_j, \qquad \forall i\in I,\ j\in C.
\end{equation}
Justificativa: se $z_j=0$, então $a_{ij}=0$; evita atribuição
a centros não instalados.

\textbf{(R4) Definição da maior distância.}
\begin{equation}
D \ge d_{ij}\, a_{ij}, \qquad \forall i\in I,\ j\in C.
\end{equation}
Justificativa: como cada $a_{ij}\in\{0,1\}$, a restrição ativa
apenas a distância do centro escolhido. Logo, $D$ é pelo menos
a distância de cada cliente ao seu centro, e minimizar $D$
equivale a minimizar a maior distância de atendimento. A
restrição é linear porque $d_{ij}$ é constante.

\textbf{(R5) Domínios.}
\begin{equation}
z_j\in\{0,1\},\quad a_{ij}\in\{0,1\},\quad D\ge 0.
\end{equation}

% ------------------------------------------------
% 3.2 ITEM B
% ------------------------------------------------
\subsection{Modelo do item (b)}

\subsubsection*{Variáveis de decisão}

$z_j$ definida como no item (a) e
\begin{equation}
u_i = \begin{cases}
1, & \text{se a localidade } i \text{ está coberta},\\
0, & \text{caso contrário.}
\end{cases}
\end{equation}

Defini-se ainda
\begin{equation}
N_i(R) = \{\,j\in C : d_{ij}\le R\,\}.
\end{equation}

\subsubsection*{Função objetivo}

\begin{equation}
\boxed{\max \; \sum_{i\in I} v_i\, u_i}
\end{equation}

\subsubsection*{Restrições e justificativas}

\textbf{(R6) Número exato de centros.}
\begin{equation}
\sum_{j\in C} z_j = m.
\end{equation}

\textbf{(R7) Cobertura condicionada à existência de centro
dentro do raio.}
\begin{equation}
u_i \le \sum_{j\in N_i(R)} z_j, \qquad \forall i\in I.
\end{equation}
Justificativa: se nenhum centro for aberto dentro do raio,
o lado direito é 0 e $u_i$ é forçado a 0. Como $v_i>0$, o
solver faz $u_i=1$ sempre que houver cobertura, ativando a
variável. Note-se que a restrição é suficiente para o modelo
de maximização, pois não há incentivo a deixar $u_i=0$ quando
há cobertura disponível.

\textbf{(R8) Domínios.}
\begin{equation}
z_j\in\{0,1\},\quad u_i\in\{0,1\}.
\end{equation}

% ------------------------------------------------
% 4. IMPLEMENTAÇÃO COMPUTACIONAL
% ------------------------------------------------
\section{Implementação Computacional}

\subsection{Solver e ambiente}

Os modelos foram implementados em Python 3, com a biblioteca
\texttt{SciPy} e a função \texttt{scipy.optimize.milp}, que
utiliza o solver HiGHS. O HiGHS é um solver de código aberto,
sem restrição de licença por tamanho de instância.

Cada instância é executada com:
\begin{itemize}
    \item \texttt{time\_limit = 3600} (1 hora, conforme o
    enunciado);
    \item \texttt{mip\_rel\_gap = 0.0} (exigência de
    otimalidade exata, para obter o gap reportado);
    \item \texttt{disp = False} (sem saída verbosa no console).
\end{itemize}

\subsection{Leitura dos dados}

Cada arquivo \texttt{inst\_n\_m.txt} tem o seguinte formato:

\begin{verbatim}
n
m
R
x1 y1 v1 e1
x2 y2 v2 e2
...
xn yn vn en
\end{verbatim}

Após a leitura, as distâncias euclidianas são calculadas
para todos os pares $(i,j)$.

\subsection{Extrato dos resultados}

Ao final de cada execução, o código registra:
\begin{itemize}
    \item valor da função objetivo;
    \item centros escolhidos ($z_j=1$);
    \item status do solver (\texttt{res.status});
    \item gap de otimalidade (\texttt{res.mip\_gap});
    \item bound (\texttt{res.mip\_dual\_bound});
    \item tempo computacional medido com
    \texttt{time.perf\_counter()}.
\end{itemize}

Para o item (b), como o modelo é de maximização, a função
objetivo é reescrita como $-\sum_i v_i u_i$ (minimização),
e o valor reportado é o negativo de \texttt{res.fun}.

\subsection{Relação com os métodos vistos em aula}

O método Simplex estudado em aula resolve problemas de
programação linear contínua. Como os modelos desta questão
exigem decisões binárias (instalar ou não um centro, atribuir
ou não um cliente, cobrir ou não uma localidade), trata-se
de problemas de programação linear inteira mista. Por isso,
usa-se um solver de PLIM (HiGHS), que aplica branch-and-bound
e cortes para impor integralidade.

% ------------------------------------------------
% 5. EXPERIMENTOS
% ------------------------------------------------
\section{Experimentos Computacionais}

Foram utilizadas nove instâncias com quantidades crescentes
de localidades e centros. As colunas obrigatórias são:
tempo computacional, status do solver, gap de otimalidade e
bound. Quando a solução é ótima, tem-se gap $=0$ e bound
igual ao valor da função objetivo.

\subsection{Resultados do item (a)}

A Tabela~\ref{tab:resultado_a} apresenta os resultados
obtidos para a minimização da maior distância.

\begin{table}[H]
\centering
\caption{Resultados do item (a) --- minimização da maior distância}
\label{tab:resultado_a}
\small
\begin{tabular}{c c c l r r l r r}
\toprule
Instância & $n$ & $m$ & Centros escolhidos & $D^*$ &
Tempo (s) & Status & Gap & Bound\\
\midrule
20\_3  & 20 & 3  & 5, 12, 19 & 182,923 & --- & Ótimo & 0,00 & 182,923\\
20\_4  & 20 & 4  & 2, 5, 8, 15 & 163,686 & --- & Ótimo & 0,00 & 163,686\\
30\_4  & 30 & 4  & 3, 13, 14, 22 & 182,833 & --- & Ótimo & 0,00 & 182,833\\
40\_8  & 40 & 8  & 2, 5, 7, 19, 20, 24, 28, 40 & 125,925 & --- & Ótimo & 0,00 & 125,925\\
40\_9  & 40 & 9  & \multicolumn{6}{l}{\textit{reexecutar com solver sem restrição de licença}}\\
50\_7  & 50 & 7  & \multicolumn{6}{l}{\textit{reexecutar com solver sem restrição de licença}}\\
50\_10 & 50 & 10 & \multicolumn{6}{l}{\textit{reexecutar com solver sem restrição de licença}}\\
60\_11 & 60 & 11 & \multicolumn{6}{l}{\textit{reexecutar com solver sem restrição de licença}}\\
60\_12 & 60 & 12 & \multicolumn{6}{l}{\textit{reexecutar com solver sem restrição de licença}}\\
\bottomrule
\end{tabular}
\end{table}

\paragraph{Observação sobre a Tabela~\ref{tab:resultado_a}.}
A execução anterior das instâncias 40\_9 a 60\_12 utilizou
um solver com limitação de licença para o tamanho do problema,
o que \emph{não} caracteriza limite de tempo. Essas
instâncias devem ser reexecutadas com o HiGHS (código do
Apêndice~\ref{ap:codigo}), que não tem essa limitação, com
\texttt{time\_limit = 3600} s. A coluna de tempo deve ser
preenchida com o valor medido na execução.

\subsection{Resultados do item (b)}

A Tabela~\ref{tab:resultado_b} apresenta os resultados
obtidos para a maximização das viagens cobertas.

\begin{table}[H]
\centering
\caption{Resultados do item (b) --- maximização de viagens cobertas}
\label{tab:resultado_b}
\small
\begin{tabular}{c c c c l r r l r r}
\toprule
Instância & $n$ & $m$ & $R$ & Centros escolhidos &
Viagens & Tempo (s) & Status & Gap & Bound\\
\midrule
20\_3  & 20 & 3  & 110 & 6, 10, 14 & 72 & --- & Ótimo & 0,00 & 72\\
20\_4  & 20 & 4  & 98  & 2, 8, 15, 18 & 86 & --- & Ótimo & 0,00 & 86\\
30\_4  & 30 & 4  & 110 & 3, 13, 24, 26 & 140 & --- & Ótimo & 0,00 & 140\\
40\_8  & 40 & 8  & 76  & 2, 10, 22, 24, 27, 35, 36, 38 & 187 & --- & Ótimo & 0,00 & 187\\
40\_9  & 40 & 9  & 87  & 3, 4, 5, 11, 16, 31, 35, 36, 40 & 222 & --- & Ótimo & 0,00 & 222\\
50\_7  & 50 & 7  & 88  & 3, 4, 11, 13, 35, 43, 50 & 235 & --- & Ótimo & 0,00 & 235\\
50\_10 & 50 & 10 & 62  & 5, 6, 8, 17, 19, 22, 26, 38, 40, 41 & 224 & --- & Ótimo & 0,00 & 224\\
60\_11 & 60 & 11 & 81  & 5, 7, 15, 22, 30, 34, 40, 43, 49, 50, 55 & 290 & --- & Ótimo & 0,00 & 290\\
60\_12 & 60 & 12 & 70  & 9, 11, 13, 17, 27, 36, 45, 46, 48, 50, 51, 59 & 331 & --- & Ótimo & 0,00 & 331\\
\bottomrule
\end{tabular}
\end{table}

\paragraph{Observação sobre a Tabela~\ref{tab:resultado_b}.}
As colunas de tempo devem ser preenchidas com o valor medido
pelo código do Apêndice~\ref{ap:codigo}. Nos casos em que o
solver retornar status ``Ótimo'', o gap é 0 e o bound
coincide com o valor da função objetivo. Caso alguma
instância não seja resolvida até a otimalidade em 3600 s,
deve-se reportar o melhor valor encontrado, o bound
(\texttt{res.mip\_dual\_bound}) e o gap
(\texttt{res.mip\_gap}).

\subsection{Análise comparativa}

No item (a), a prioridade é reduzir a maior distância de
atendimento, o que leva a uma distribuição espacial
equilibrada dos centros. No item (b), a prioridade é
maximizar as viagens cobertas, levando em conta os valores
$v_i$ e o raio $R$. Por isso, as localidades escolhidas
podem ser diferentes entre os dois itens.

Observa-se também que aumentar o número de centros tende a
reduzir $D^*$ no item (a) e aumentar as viagens cobertas no
item (b), mas o efeito depende da distribuição espacial das
localidades e das posições elegíveis em $C$.

% ------------------------------------------------
% 6. CONCLUSÃO
% ------------------------------------------------
\section{Conclusão}

O trabalho apresentou a formulação e a resolução
computacional de dois problemas de localização de centros
de distribuição.

No item (a), foi desenvolvido um modelo de minimização da
maior distância de atendimento, com restrições que garantem
exatamente $m$ centros, atribuição única dos clientes e
proibição de atribuição a centros não instalados.

No item (b), foi desenvolvido um modelo de maximização das
viagens cobertas, com variáveis binárias de cobertura e
restrições que exigem, para cada cliente, a existência de
pelo menos um centro instalado dentro do raio $R$.

Os experimentos com as nove instâncias mostraram como o
número de centros, a distribuição espacial das localidades
e o raio de cobertura influenciam os resultados. A
implementação com solver de PLIM (HiGHS) garante que as
decisões de instalação, atribuição e cobertura respeitem
a integralidade.

Por fim, a modelagem matemática mostrou-se fundamental para
representar corretamente problemas reais de localização, e
a utilização de métodos de otimização inteira permitiu obter
soluções comprovadamente ótimas para as instâncias
estudadas.

% ------------------------------------------------
% REFERÊNCIAS
% ------------------------------------------------
\section{Referências}

\begin{itemize}
    \item Material de aula: Forma Padrão.
    \item Material de aula: Gabarito da Revisão para a Prova.
    \item Arquivos de dados das instâncias disponibilizadas
    para o Trabalho Final 1.
    \item Documentação da biblioteca SciPy:
    \url{https://docs.scipy.org/doc/scipy/reference/optimize.milp.html}.
    \item HiGHS --- \url{https://highs.dev/}.
\end{itemize}

% ------------------------------------------------
% APÊNDICE
% ------------------------------------------------
\appendix

\section{Código-fonte da implementação}
\label{ap:codigo}

https://bito.li/NDlmb

\end{document}